\documentclass[11pt]{article}
% Shared preamble for the qurl-theory reports.
% Build either report from this folder with:  latexmk -pdf <file>.tex
% (or: pdflatex <file>; bibtex <file>; pdflatex <file>; pdflatex <file>)

\usepackage[T1]{fontenc}
\usepackage{lmodern}
\usepackage[margin=1in]{geometry}
\usepackage{microtype}
\usepackage{amsmath,amssymb,amsthm,mathtools}
\usepackage{booktabs,tabularx,array}
\usepackage{enumitem}
\usepackage[dvipsnames]{xcolor}
\usepackage[authoryear,round]{natbib}
\usepackage{xurl}  % lets long URLs break anywhere (loads url)

% Allow a line break (without an inserted hyphen) after each \_ in text, so long file names can wrap.
\makeatletter
\DeclareRobustCommand{\_}{\ifmmode\nfss@text{\textunderscore}\else\textunderscore\allowbreak\fi}
\makeatother
\usepackage[colorlinks=true,linkcolor=NavyBlue,citecolor=ForestGreen,urlcolor=RoyalBlue]{hyperref}
\usepackage[capitalise,noabbrev]{cleveref}

% ---- theorem-like environments (defined after cleveref so shared counters get the right names)
\theoremstyle{plain}
\newtheorem{theorem}{Theorem}
\newtheorem{proposition}[theorem]{Proposition}
\newtheorem{lemma}{Lemma}
\newtheorem{corollary}{Corollary}
\theoremstyle{definition}
\newtheorem{definition}{Definition}
\newtheorem*{assumptionU}{Assumption U}
\theoremstyle{remark}
\newtheorem{remark}{Remark}

% ---- status tags (robust, so they can appear in section titles and the table of contents)
\DeclareRobustCommand{\statustag}[2]{\textsf{\footnotesize\textcolor{#1}{[#2]}}}
\DeclareRobustCommand{\tagchecked}{\statustag{ForestGreen}{checked}}
\DeclareRobustCommand{\tagproved}{\statustag{RoyalBlue}{proved}}
\DeclareRobustCommand{\tagderived}{\statustag{RoyalBlue}{derived}}
\DeclareRobustCommand{\tagconj}{\statustag{BurntOrange}{conjecture}}
\DeclareRobustCommand{\tagopen}{\statustag{BrickRed}{open}}

% ---- research-problem headings (numbered P1, P2, ... by hand, listed in the table of contents)
\newcommand{\researchproblem}[2]{%
  \subsubsection*{#1.\ #2}\addcontentsline{toc}{subsubsection}{#1.\ #2}}

% ---- math and text macros
\newcommand{\E}{\mathbb{E}}
\newcommand{\R}{\mathbb{R}}
\newcommand{\N}{\mathbb{N}}
\newcommand{\Z}{\mathbb{Z}}
\newcommand{\Unif}{\mathcal{U}}
\DeclareMathOperator{\Var}{Var}
\DeclareMathOperator{\Corr}{Corr}
\DeclareMathOperator{\clip}{clip}
\DeclareMathOperator{\diag}{diag}
\DeclareMathOperator{\sgn}{sign}
\DeclareMathOperator{\fracpart}{frac}
\DeclareMathOperator{\rms}{rms}
\newcommand{\KL}{\mathrm{KL}}
\newcommand{\TV}{\mathrm{TV}}
\newcommand{\script}[1]{\texttt{sanity\_checks/\allowbreak#1}}
\newcolumntype{L}{>{\raggedright\arraybackslash}X}
\setlist{itemsep=2pt,topsep=4pt}

% ---- layout: avoid overfull lines from unbreakable math/citations; ragged-right bibliography
\setlength{\emergencystretch}{1.5em}
\renewcommand{\bibfont}{\raggedright}


\title{P1: UAQ and Update Tracking\\in Quantized-Rollout Reinforcement Learning}
\author{Theory note for research problem P1 (draft)}
\date{7 November 2026; reorganized 11 November 2026; experiments split off 12 November 2026; E2's results added 15 November 2026; statements corrected after review 18 November 2026}

\begin{document}
\maketitle

\paragraph{Scope.} This is the theory note for research problem P1 in the companion report \texttt{research\_problems.tex}. Its subject is QuRL's Update-Aware Quantization (UAQ) and its ``weight update problem'' (Sec.~4.3 and App.~A of \citealp{li2026qurl}). On 12 November 2026 the experiments moved to the companion experiments report \texttt{p1\_experiments.tex}, which tests the results below.

\section*{In short}\addcontentsline{toc}{section}{In short}
\begin{itemize}
\item UAQ leaves the quantized rollout model exactly unchanged (\cref{thm:invariance}). Where UAQ changes a quantized model at all, it does so through a different rounding pattern, never by an $s$-fold error reduction (\cref{rem:scope}).
\item UAQ's whole effect on training is a change of optimizer settings, chiefly per-layer learning rates: $\times s$ on Q/K/V and $\div s$ on the norm gain under AdamW, with $\varepsilon$, weight decay and the clipping norm matched (\cref{thm:reparam}). QuRL + UAQ is therefore identical to QuRL with these settings (\cref{cor:pipeline}), and the ``$s^2$ improvement'' is $s$ under Adam (\cref{cor:ratio}).
\item Re-quantized master weights do not freeze; updates are captured in expectation (\cref{prop:lag,prop:capture,prop:layer}). The real issue is that the rounding error persists across steps (\cref{prop:persistence}). Error feedback is the scheme whose cumulative error stays bounded (\cref{prop:schemes}).
\item \textbf{Evidence.} The companion experiments report (\texttt{p1\_experiments.tex}) tests these results in float64, on vLLM's kernels, in an exact training loop and in pre-registered RL runs. So far:
  \begin{itemize}
  \item \cref{thm:invariance,thm:reparam} hold bit-for-bit on the real stack;
  \item on the master weights of an instrumented RL run (E2), \cref{prop:persistence,prop:schemes} hold, and so does \cref{prop:capture} from the second update on (\cref{rem:bf16});
  \item nothing contradicts any result where its assumptions hold.
  \end{itemize}
  The report's Section~4 goes through the evidence claim by claim.
\end{itemize}

\tableofcontents

% =====================================================================================
\section{Background: how the problem arises}\label{sec:background}

\subsection{Quantized rollouts}
\begin{itemize}
\item \textbf{RL for LLMs} alternates generating responses (rollouts) and training on them. Generation dominates the step time.
\item \textbf{QuRL} generates with a W8A8-quantized copy of the actor (rollout policy $\mu_\theta$) and trains full-precision (FP) master weights (policy $\pi_\theta$). After every update the master weights are re-quantized and sent to the rollout engine.
\item \textbf{The cost} is a train--rollout mismatch, $\mu\neq\pi$, which makes training off-policy. QuRL corrects it with truncated importance weights $\min(\pi_k/\mu_k,C)$ (TIS, \citealp{yao2025flashrl}) and a widened clip bound (ACR).
\item \textbf{Step by step.} The QuRL explainer (\texttt{qurl\_explained.tex}) gives the whole algorithm (its Algorithm~1): re-quantization is line~3, the 8-bit rollout line~4, TIS line~7, ACR line~10 and UAQ line~1.
\end{itemize}

\subsection{Absmax quantization and the grid step}
\begin{itemize}
\item \textbf{Quantizer.} $Q(x)=\alpha\,r(x/\alpha)$ with $\alpha=\max_{\text{group}}|x|/M$: one scale per output channel for weights, per token for activations. For INT8, $M=127$ and the grid step is $\Delta=\alpha$.
\item \textbf{Scales follow the tensor.} Multiply a group by $c>0$ and its scale is multiplied by $c$.
\item \textbf{Error.} Round-to-nearest (RTN) is off by at most $\Delta/2$.
\end{itemize}

\subsection{QuRL's two claims}\label{sec:claims}
\begin{itemize}
\item \textbf{(i) The weight update problem.} Per-step updates (about $10^{-7}$ to $10^{-6}$) are far below $\Delta$, so re-quantization ``masks nearly all weight updates, effectively freezing the quantized model'' (Sec.~4.3).
\item \textbf{(ii) UAQ fixes it.} Once before RL, set $W_{Q,K,V}\to W/s$ and the preceding RMSNorm gain $\gamma\to s\gamma$, with $s=1.5$. QuRL applies $s$ ``column-wise to $W$ and row-wise to $X$'' but reports only scalar values, so $s$ is one scalar. The FP model is unchanged. The paper says this ``reduces quantization error by a factor of $s$ while amplifying weight updates by the same factor'', an ``$s^2$ improvement'' (Sec.~4.3). The explainer's Section~3.7 summarizes both claims as the paper makes them.
\item \textbf{Evidence offered.} Mismatch curves with and without UAQ (Fig.~4b); faster early learning with UAQ (Fig.~7); an ablation against a larger global learning rate, one number per configuration (Table~4). The reported gains from UAQ are small single runs: $+0.62$ and $+0.15$ Avg@32 points on DAPO with INT8 and FP8 (Table~2), and $+0.69$ on DeepScaleR's five-benchmark average (Table~3). On GSM8K the paper disables UAQ ``due to the relatively high learning rate ($10^{-5}$)''.
\end{itemize}

\subsection{Why the claims are doubtful}
\begin{itemize}
\item \textbf{Scales move with the tensor.} Dividing $W$ by $s$ divides its scale by $s$, so its relative error is unchanged. Multiplying $X$ by $s$ multiplies its error by $s$. In the product the two cancel.
\item \textbf{Re-quantization is not masking.} With FP master weights, a small update accumulates until it crosses a rounding boundary; at any step it flips a code with some probability.
\item \textbf{Rescaling a parameter changes its effective learning rate.} Training $W/s$ instead of $W$ changes the step size $W$ sees (the parameterization view of \citealp{yang2021tp4,yang2022tp5}).
\end{itemize}

\subsection{Notation}\label{sec:notation}
\begin{itemize}
\item \textbf{Linear layer.} $Y=XW^\top$ with $X\in\R^{T\times d}$ (tokens $\times$ input channels) and $W\in\R^{m\times d}$. The rollout engine computes $\hat Y=Q_a(X)\,Q_w(W)^\top$.
\item \textbf{UAQ coordinates.} RMSNorm feeds Q/K/V: $X=\gamma\odot n(x)$ with $n(x)=x/\rms(x)$. UAQ trains $(W',\gamma')=(W/s,\ s\gamma)$. The map back, $\phi_s(W',\gamma')=(sW',\gamma'/s)$, leaves the FP forward pass unchanged. Scaling the MLP inputs through the second norm works the same way.
\item \textbf{QuRL step $k$.} Rollouts come from $\mu_k=\mu_{\theta_k}$; the proximal policy is $\pi_k=\pi_{\theta_k}$. The loss is the decoupled-PPO surrogate \citep{hilton2022batch} with TIS weights and ACR (or PPO) clip bounds, plus an optional KL term to $\pi_{\mathrm{ref}}$. It depends on the parameters only through the policies and the sampled data. Updates are minibatch steps of an optimizer $\mathcal O$, optionally with global-norm gradient clipping.
\end{itemize}

% =====================================================================================
\section{Problem statements}\label{sec:problems}
\begin{itemize}
\item \textbf{Q1. Does UAQ change the rollout model?} For which quantizers is $Q_a(sX)\,Q_w(W/s)^\top=Q_a(X)\,Q_w(W)^\top$? $\to$ \cref{thm:invariance}.
\item \textbf{Q2. What does UAQ change in training?} Is training in UAQ coordinates the same as training the original parameters under other optimizer settings, and which ones? $\to$ \cref{thm:reparam}, \cref{cor:ratio,cor:pipeline}.
\item \textbf{Q3. Do re-quantized master weights freeze?} How much of a sub-grid update reaches the rollout model? $\to$ \cref{prop:lag,prop:capture,prop:layer}.
\item \textbf{Q4. What is the real issue, and what should ``update-aware'' mean?} How does the rounding error evolve during training, and which rollout-quantization scheme controls its effect? $\to$ \cref{prop:persistence,prop:schemes}, \cref{cor:persistence-time,cor:partial-sums}.
\item \textbf{Q5. Do the answers hold on a real RL stack?} On vLLM's kernels, in an exact training loop, and in distribution over seeds in full RL runs (hypotheses H1--H3). $\to$ the experiments report (\texttt{p1\_experiments.tex}).
\end{itemize}

% =====================================================================================
\section{Theory: statements and proofs}\label{sec:theory}

\subsection{UAQ leaves the quantized forward pass unchanged (Q1)}\label{sec:invariance}

\begin{definition}[Homogeneous group quantizer]
Partition a tensor's entries into groups $g\in\mathcal P$. $Q$ is \emph{homogeneous on $\mathcal P$} if $Q(c\odot X)=c\odot Q(X)$ for every positive $c$ that is constant on each group.
\end{definition}

A single-level absmax quantizer is $Q(X)_i=\alpha_g\,r(X_i/\alpha_g)$ for $i\in g$, with $\alpha_g=\max_{j\in g}|X_j|/M$ and $r$ a fixed rounding onto a grid with maximum $M$ (INT-$b$: $M=2^{b-1}-1$; FP8-E4M3: 448; FP4-E2M1: 6).

\begin{lemma}\label{lem:homogeneous}
Single-level absmax quantizers are homogeneous on their groups. Multi-level quantizers in which every scale is an absmax or a rounding of a \emph{ratio} of absmaxes are homogeneous on the groups of their coarsest scale, not on their finer blocks. For example, NVFP4 (E4M3 block scales under a per-tensor FP32 absmax scale, \citealp{nvidia2025nvfp4}) is homogeneous on the whole tensor.
\end{lemma}
\begin{proof}
If $c=c_g$ on group $g$, then $\alpha_g(c\odot X)=c_g\alpha_g(X)$, so $X_i/\alpha_g$ is unchanged and $Q(c\odot X)_i=c_gQ(X)_i$. In the multi-level case, with $c$ constant on each group of the coarsest scale, every absmax in such a group scales by the same factor, so the rounded ratios are unchanged. A factor that differed between finer blocks would change them.
\end{proof}

\begin{theorem}\label{thm:invariance}
Let $X'=X\diag(s)$ and $W'=W\diag(s)^{-1}$ with $s\in\R^d_{>0}$ constant on the input-channel footprint of every group on which $Q_a$ or $Q_w$ is homogeneous (\cref{lem:homogeneous}). Then $Q_a(X')\,Q_w(W')^\top=Q_a(X)\,Q_w(W)^\top$. This holds:
\begin{itemize}
\item for QuRL's W8A8 (per-token activations, per-output-channel weights), for every \textbf{scalar} $s>0$;
\item for single-level block quantization, for any $s$ constant on blocks;
\item for two-level scaling such as NVFP4, for scalar $s$ only.
\end{itemize}
\end{theorem}
\begin{proof}
\Cref{lem:homogeneous} twice: $Q_a(X')=Q_a(X)\diag(s)$ and $Q_w(W')=Q_w(W)\diag(s)^{-1}$.
\end{proof}

\begin{corollary}\label{cor:rollout}
If every quantized layer touched by UAQ satisfies \cref{thm:invariance}, then $\pi_{\phi_s^{-1}(\theta)}=\pi_\theta$ and $\mu_{\phi_s^{-1}(\theta)}=\mu_\theta$. With the same random seeds, rollouts, rewards, advantages, ratios $\pi/\mu$ and every divergence between $\mu$ and $\pi$ are identical before and after UAQ.
\end{corollary}

\begin{remark}[The paper's error formula]
``$\hat\theta-\theta\propto|\theta|/(s2^b)$'' is the absolute weight error in rescaled units: $|Q(W/s)-W/s|=|Q(W)-W|/s$, while $|Q(sX)-sX|=s|Q(X)-X|$. The two cancel exactly in $\hat Y$.
\end{remark}

\begin{remark}[Where \cref{thm:invariance} does not apply]\label{rem:scope}
Relative output error of one quantized layer, from \script{uaq\_edge\_cases.py}:
\begin{center}\small
\begin{tabularx}{\linewidth}{@{}L l l L@{}}
\toprule
Quantizer & $s=1.5$ & $s=2$ & Why\\
\midrule
INT8, per-channel weights, per-token activations & unchanged ($3\times10^{-16}$) & unchanged & \Cref{thm:invariance}\\
NVFP4-like (E4M3 blocks under FP32 tensor scale) & unchanged ($2\times10^{-16}$) & unchanged & \Cref{lem:homogeneous}\\
MXFP4 (power-of-two block scales, \citealp{rouhani2023mx}) & 18.42\% $\to$ 18.70\% & unchanged & E8M0 is homogeneous only for $2^k$\\
INT8, static activation scale, not recalibrated & 7.65\% $\to$ 9.28\% & 18.7\% & becomes a clipping-range change\\
FP8 cast without scale & 4.12\% $\to$ 4.57\% & 5.06\% & subnormals have a fixed grid\\
INT8 of BF16-stored weights & \parbox[t]{9em}{\raggedright 3.7\% of codes re-rounded; 4.822\% $\to$ 4.829\%} & unchanged & BF16 rounding of $W/s$\\
\bottomrule
\end{tabularx}
\end{center}
Where UAQ changes the quantized model at all, it does so through a different rounding pattern or a calibration side-effect, never by an $s$-fold error reduction.
\end{remark}

\subsection{UAQ is an optimizer reparameterization (Q2)}\label{sec:reparam}

\begin{lemma}[Gradients]\label{lem:gradients}
With $\mathcal L'=\mathcal L\circ\phi_s$, the chain rule gives $\nabla_{W'}\mathcal L'=s\,\nabla_W\mathcal L$ and $\nabla_{\gamma'}\mathcal L'=s^{-1}\nabla_\gamma\mathcal L$; all other gradients are unchanged.
\end{lemma}

\begin{theorem}\label{thm:reparam}
Run optimizer $\mathcal O$ on $\theta'=(W',\gamma',\mathrm{rest})$ with learning rate $\eta$, Adam constant $\varepsilon$ and decoupled weight decay $\lambda$ \citep{kingma2015adam,loshchilov2019adamw}. Run the same $\mathcal O$ on $\theta=\phi_s(\theta')$ with these per-group settings:
\begin{center}\small
\begin{tabularx}{\linewidth}{@{}L l l l l@{}}
\toprule
Optimizer & lr ($W$) & lr ($\gamma$) & $\varepsilon$ ($W$ / $\gamma$) & decay on $W$ ($\gamma$)\\
\midrule
(a) SGD, heavy-ball, Nesterov & $s^2\eta$ & $\eta/s^2$ & --- & $\lambda/s^2$ ($s^2\lambda$)\\
(b) Adam, AdamW, RMSProp, Adagrad & $s\eta$ & $\eta/s$ & $\varepsilon/s$ / $s\varepsilon$ & $\lambda/s$ ($s\lambda$)\\
(b$'$) sign-SGD, Lion & $s\eta$ & $\eta/s$ & --- & $\lambda/s$ ($s\lambda$)\\
\bottomrule
\end{tabularx}
\end{center}
(c) \textbf{Clipping.} With global-norm clipping, compute the norm on $(s\,g_W,\ g_\gamma/s,\ g_{\mathrm{rest}})$, that is, in UAQ coordinates.

Then $\theta_k=\phi_s(\theta'_k)$ for every step $k$.
\end{theorem}
\begin{proof}
Induction on $k$. Hypothesis $H_k$: $\theta_k=\phi_s(\theta'_k)$, and the optimizer states correspond ($m'=s\,m$, $v'=s^2v$ for $W$; $m'=m/s$, $v'=v/s^2$ for $\gamma$; momentum likewise).
\begin{itemize}
\item \emph{Base.} $H_0$ holds by construction, with zero states.
\item \emph{Same data, matching gradients.} By \cref{cor:rollout} both runs see identical data; by \cref{lem:gradients} the gradients are $(s\,g_W,\ g_\gamma/s)$.
\item \emph{Same clip.} Under (c) the clip factor is the same number in both runs.
\item \emph{States.} Moment updates are linear in $g$ and $g^2$, so the correspondence carries to $k+1$.
\item \emph{Adam step on $W$:}
\[
\Delta W'=-\eta\,\frac{s\,\hat m}{s\sqrt{\hat v}+\varepsilon}-\eta\lambda W' \;\Longrightarrow\; \Delta W=s\,\Delta W'=-s\eta\,\frac{\hat m}{\sqrt{\hat v}+\varepsilon/s}-\eta\lambda W,
\]
which is AdamW with lr $s\eta$, constant $\varepsilon/s$ and decay $\lambda/s$, since $(s\eta)(\lambda/s)=\eta\lambda$. On $\gamma$: $\Delta\gamma=\Delta\gamma'/s=-(\eta/s)\,\hat m/(\sqrt{\hat v}+s\varepsilon)$.
\item \emph{SGD:} $\Delta W=s\cdot(-\eta s g_W)=-s^2\eta g_W$ and $\Delta\gamma=-(\eta/s^2)g_\gamma$. \emph{Sign-based optimizers} ignore positive per-group gradient scales.\qedhere
\end{itemize}
\end{proof}

\begin{corollary}[The ``$s^2$ improvement'']\label{cor:ratio}
The update-to-grid ratio $\nu_{ij}=|\Delta W_{ij}|/\Delta_i(W)$ ($\Delta_i$: row $i$'s grid step) is invariant under $\phi_s$. So UAQ's $\nu$ equals the matched run's. At the first step from a shared starting point, that is $s^p$ times the plain run's: $p=2$ for SGD, and $p=1$ for Adam-type optimizers as $\varepsilon\to0$. After that the runs diverge. VeRL trains with AdamW, so the first-step gain is $s$, not $s^2$. Raising a layer's $\nu$ and raising its learning rate are the same operation.
\end{corollary}

\begin{corollary}[The whole pipeline]\label{cor:pipeline}
\Cref{cor:rollout} at every iterate plus \cref{thm:reparam} make QuRL + UAQ at learning rate $\eta$ \emph{identical} to QuRL without UAQ under \cref{thm:reparam}'s settings: rollouts, rewards, TIS weights, PPO/ACR clip decisions, the KL-to-reference term and every parameter. This holds for any rollout correction that depends only on policy values (TIS, ACR, MIS, sequence masks). \textbf{UAQ's entire effect on training is an effect of per-group optimizer settings}, chiefly the learning rates.
\end{corollary}

\begin{remark}[Finite precision]\label{rem:practice}
Bit-exact equivalence needs exact rescaling and deterministic kernels. At $s=2$ every rescaling is exact in BF16 and FP32; at $s=1.5$, $W/s$ is rounded, and RL amplifies the difference as it amplifies a change of seed. That UAQ and the matched settings are then \emph{equal in distribution over seeds} is a hypothesis, not a consequence of \cref{cor:pipeline}: deterministic rounding differences could also change how often runs enter a different training regime. It is hypothesis H1 of the experiments report (its Section~2.4).
\end{remark}

\subsection{Re-quantized master weights track updates (Q3)}\label{sec:tracking}

\begin{assumptionU}[Uniform phase]
For a coordinate with grid step $\Delta$, write $w/\Delta=c+\varphi$, $c\in\Z$, $\varphi\in[-\tfrac12,\tfrac12)$; RTN gives $Q(w)=c\Delta$ and error $e=-\Delta\varphi$. Assume $\varphi\sim\Unif[-\tfrac12,\tfrac12)$, independent of the update (over several steps, of the cumulative displacement) and across coordinates. In RL the path depends on rollouts from the quantized weights, so the independence is itself an assumption. E2 finds phase and one-step update uncorrelated ($|r|\le6\times10^{-4}$).
\end{assumptionU}

\begin{lemma}[When U holds]\label{lem:phase}
If $w$ has density $p$ with characteristic function $\hat p$, the phase density $f$ satisfies $\sup_\varphi|f(\varphi)-1|\le\sum_{m\ne0}|\hat p(2\pi m/\Delta)|$. For a Gaussian (or mixture) with standard deviation at least $\sigma$, the bound is $2\sum_{m\ge1}e^{-2\pi^2m^2\sigma^2/\Delta^2}<10^{-8}$ already at $\sigma=\Delta$ \citep{widrow2008quantization,sripad1977necessary}. Per-row absmax INT8 has $\sigma/\Delta\approx127/(\max/\rms)$. On the Qwen2.5 checkpoints of 0.5B, 1.5B and 7B parameters, $\max/\rms$ has a median of 3.6--4.5 per projection type, so $\sigma/\Delta$ is 28--36, and no row has $\sigma/\Delta<1$ (experiments report, E2). U fails if weights sit on grid points, e.g.\ after QAT on the same grid.
\end{lemma}
\begin{proof}
$f(\varphi)=\sum_k\Delta\,p(\Delta(k+\varphi))$ is 1-periodic with Fourier coefficients $\hat p(2\pi m/\Delta)$ and $\hat p(0)=1$ (Poisson summation).
\end{proof}

\begin{remark}[BF16 storage]\label{rem:bf16}
\Cref{lem:phase} assumes a density. Weights stored in BF16 have none: they sit on a lattice whose spacing, in grid units, is $|c|/256$ to $|c|/128$ for code $c$.
\begin{itemize}
\item U then holds at the resolution of a few hundredths of a grid step, but not at the scale of a single update.
\item Next to a rounding boundary, the stored weights sit either on it (exact ties, 0.07--0.26\% of the weights on the three checkpoints above) or at discrete distances that are often larger than one update. So fewer weights lie within one update of a boundary than U assumes.
\item Training with FP32 master weights moves the weights off the lattice within a few steps.
\end{itemize}
In E2 (experiments report), \cref{prop:capture}'s capture slope is 0.93 for the first update from the BF16 checkpoint, 0.97 for the second and 0.99 from step 64 on. A simulation reproduces the first two values from BF16 storage alone.
\end{remark}

\begin{proposition}[Bounded lag; no assumption]\label{prop:lag}
For RTN on a fixed grid without saturation, $\|Q(\theta_k)-\theta_k\|_\infty\le\Delta/2$ for all $k$, so $\|(Q(\theta_k)-Q(\theta_0))-(\theta_k-\theta_0)\|_\infty\le\Delta$. The quantized trajectory never falls more than one grid step behind. While the total drift stays below one grid step, as in QuRL's regime, this alone does not exclude freezing; \cref{prop:capture} does.
\end{proposition}

\begin{proposition}[One-step capture, under U]\label{prop:capture}
Write $|d|/\Delta=n_0+a$, $n_0\in\N$, $a\in[0,1)$. Then $Q(w+d)-Q(w)=\sgn(d)\,\Delta\,(n_0+B)$ with $B\sim\mathrm{Bernoulli}(a)$. Hence $\E[Q(w+d)-Q(w)]=d$ (captured in expectation), $\Var=\Delta^2a(1-a)$ ($=\Delta|d|-d^2$ for $|d|<\Delta$), and for $|d|<\Delta$ the code flips with probability $|d|/\Delta$.
\end{proposition}
\begin{proof}
$(w+d)/\Delta=c+n_0+(\varphi+a)$ rounds to $c+n_0+\mathbf 1[\varphi\ge\tfrac12-a]$, an indicator with probability $a$.
\end{proof}

\begin{proposition}[Layer-level tracking]\label{prop:layer}
For a weight row with independent U-phases, update $d$ and input $x$: $\E[(Q(w+d)-Q(w))\cdot x]=d\cdot x$ and $\Var=\sum_jx_j^2\Delta_j^2a_j(1-a_j)$. With $|d_j|=\epsilon\Delta$:
\begin{itemize}
\item gradient-aligned update ($d_j=\epsilon\Delta\,\sgn(x_j)$): $\mathrm{SNR}=\sqrt{\epsilon/(1-\epsilon)}\,\|x\|_1/\|x\|_2\approx\sqrt{2\epsilon n/\pi}$ for Gaussian $x$;
\item random-sign update: $\mathrm{SNR}\approx\sqrt{\epsilon/(1-\epsilon)}$, independent of $n$.
\end{itemize}
\end{proposition}
To first order the same holds in policy space with $x$ replaced by $\nabla_\theta\log\pi(a\mid c)$. Gradient updates are aligned with the scores of the tokens that produced them, so the rollout policy tracks the \emph{intended} changes well.

\subsection{The real issue: persistent error (Q4)}\label{sec:persistence}

\begin{proposition}[Stationary but persistent]\label{prop:persistence}
Under U with a fixed grid:
\begin{enumerate}[label=(\roman*)]
\item $e_k=Q(\theta_k)-\theta_k$ has independent $\Unif[-\Delta/2,\Delta/2]$ coordinates for every $k$, so $\E\|e_k\|^2=n\Delta^2/12$ throughout training;
\item for a coordinate with cumulative displacement $D=\theta_k-\theta_0$ and $a=\fracpart(|D|/\Delta)$, $\Corr(e_k,e_0)=1-6a(1-a)$.
\end{enumerate}
\end{proposition}
\begin{proof}
(i) A uniform phase shifted by an independent amount, modulo 1, is uniform. (ii) With $\varphi'=\varphi+a-\mathbf 1[\varphi\ge\tfrac12-a]$, $\E[\varphi\varphi']=\tfrac1{12}-\int_{1/2-a}^{1/2}\varphi\,d\varphi=\tfrac1{12}-\tfrac{a(1-a)}2$; divide by $\E\varphi^2=\tfrac1{12}$.
\end{proof}

\paragraph{Consequence.} The size of the weight quantization error does not grow during training. A growing train--rollout mismatch (QuRL Fig.~3a) must come from how the errors act on the output, or from changing contexts:
\begin{itemize}
\item a growing grid (outliers);
\item growing sensitivity (Fisher, Jacobians), including sensitive directions that come to align with the persistent error;
\item growing activation error.
\end{itemize} This is research problem P4. In E2's run the weight error kept its size and the weight grids moved by at most 1.3\% for 99\% of rows, while the mismatch on a fixed set of probe tokens rose 3.5-fold. On fixed contexts, that leaves sensitivity and activation error.

\begin{corollary}[Persistence time]\label{cor:persistence-time}
The correlation falls to $\tfrac12$ at $a_{1/2}=(1-\sqrt{2/3})/2\approx0.092$: after $\approx0.092\,\Delta/|d|$ steps of coherent drift, or $\approx\tfrac\pi2(0.092\,\Delta/|d|)^2$ steps of random-walk drift.
\end{corollary}
Reading QuRL's App.~A curves as cumulative relative norms (update $\approx0.002$, quantization error $\approx0.0097$ after ${\sim}900$ steps) gives $a\approx0.06$ and $\Corr(e_{900},e_0)\approx0.66$ ($\approx0.32$ if the plotted values are squared ratios). Either way, the error at the end of training is still substantially correlated with the error at the start. E2 measures this directly: after 435 steps at $\eta=10^{-6}$ the correlation is 0.691, against 0.692 predicted from the measured displacements.

\begin{proposition}[Rollout-quantization schemes]\label{prop:schemes}
For any master-weight path $\theta_k$, SR's errors are martingale differences and error feedback's errors sum to $r_0-r_K$, with no assumption. The variances and correlations below assume uniform phases (RTN, delta) or independent, uniform residuals (SR, error feedback). A path that stays on a grid point has zero error under both.
\begin{center}\small
\begin{tabularx}{\linewidth}{@{}l L l l L L@{}}
\toprule
Scheme & Rollout weights $q_k$ & Error var. & Lag-1 corr. & $\sum_{k\le K}e_k$ & Extra cost\\
\midrule
RTN & $Q(\theta_k)$ & $\Delta^2/12$ & $\approx1$ & $\approx Ke_0$ while drift $\ll\Delta$ & none\\
Delta & $Q(\theta_0)+Q_f(\theta_k-\theta_0)$ & $\approx\Delta^2/12$ & $\approx1$ & $\approx Ke_0$ always & second GEMM or sparse delta\\
SR, fresh & $\Delta\lfloor\theta_k/\Delta+\xi_k\rfloor$ & $\Delta^2/6$ & 0 & martingale, $O(\sqrt K\Delta)$ & random bits\\
Error feedback & $Q(\theta_k+r_{k-1})$, $r_k=\theta_k+r_{k-1}-q_k$ & $\Delta^2/6$ & $-1/2$ & $r_0-r_K$; at most $\Delta/2$ for all $K$ if $r_0=0$ & one FP residual per weight\\
\bottomrule
\end{tabularx}
\end{center}
\end{proposition}
\begin{proof}
SR \citep{gupta2015limited}: $\E[e_k\mid\mathcal F_{k-1},\theta_k]=0$ exactly (no U needed), so partial sums are a martingale; variance $\Delta^2a(1-a)$ averages to $\Delta^2/6$ for a uniform fractional part $a$. Error feedback \citep{gray1990spectra,seide2014onebit,stich2018sparsified,karimireddy2019ef}: $|r_k|\le\Delta/2$ and $e_k=r_{k-1}-r_k$, so the sum telescopes for \emph{any} path; with roughly uniform independent residuals the variance is $\Delta^2/6$ and the lag-1 correlation $-1/2$. Delta: $e_k=e_0+e_{f,k}$ with $|e_{f,k}|\le\Delta_f/2$.
\end{proof}

\begin{corollary}[Why partial sums matter]\label{cor:partial-sums}
Suppose step $k$'s expected update carries a first-order mismatch bias $b_k=B_ke_k$ with $\|B_k\|\le\beta$. Let $V_K=\sum_{k<K}\|B_{k+1}-B_k\|$ be $B$'s total variation, which is at most $(K-1)\lambda$ if $\|B_{k+1}-B_k\|\le\lambda$. The accumulated bias is up to ${\sim}K\beta\|e\|$ for RTN and delta, and $O(\sqrt K\beta\|e\|)$ for SR. For error feedback (Abel summation, $r_0=0$),
\[
\Bigl\|\sum_{k=1}^K B_k(r_{k-1}-r_k)\Bigr\|\le\bigl(\beta+V_K\bigr)\max_k\|r_k\|,
\]
which is bounded whenever $B$'s total variation is.
\end{corollary}
So ``update-aware'' should mean \textbf{bounded cumulative error}, not a larger update-to-grid ratio. Turning this into a convergence theorem is research problem P6.

\subsection{What this means for QuRL's evidence}\label{sec:evidence}
These points are consistent with \cref{thm:invariance,thm:reparam}; they are not statements about the authors' code, which is not public.
\begin{itemize}
\item \textbf{Fig.~4b.} With and without UAQ, the mismatch starts at the same value (about 0.0106), as \cref{cor:rollout} predicts; the later gap is what two runs with different per-layer learning rates would show.
\item \textbf{Fig.~7 (DAPO, FP8).} QuRL + UAQ learns faster than the \emph{unquantized BF16 baseline} for ${\sim}40$ steps. Better rollout fidelity cannot beat no quantization; a higher Q/K/V learning rate can. By step 200 the two QuRL curves end level (0.278 and 0.277); each is a single run.
\item \textbf{Table~4.} The control raises the \emph{global} learning rate. The matched control (Q/K/V $\times1.5$, norm gains $\div1.5$) and BF16 + UAQ are missing, and the three $s$ settings differ by at most 2.1 Avg@32 points on 30 problems, one number each.
\item \textbf{GSM8K.} Disabling UAQ ``due to the relatively high learning rate'' is consistent with UAQ being a learning-rate change.
\item \textbf{Side hypothesis (bridge to P4).} Under AdamW, UAQ also makes the norm gains learn $s\times$ slower, which may slow outlier growth in activations and hence mismatch growth.
\end{itemize}

% =====================================================================================
\section{Open ends}\label{sec:open}
\begin{enumerate}
\item \textbf{Policy-space tracking beyond first order.} Extend \cref{prop:layer} with a Lipschitz-Hessian remainder, and add activation quantization (also recomputed from $\theta$, so also persistent).
\item \textbf{Moving grids.} Dynamic absmax rescales a row's grid whenever its largest element moves. Quantify this jitter for outlier-dominated rows and compare with frozen scales.
\item \textbf{Per-layer learning rates for RL fine-tuning.} \Cref{thm:reparam} turns UAQ into a question with nothing quantization-specific in it: which per-layer learning-rate multipliers (in the spirit of \citealp{yang2022tp5}) are best for RLVR? UAQ may be a lucky instance.
\item \textbf{Bridges.} P4: the size of the weight error is stationary, so mismatch growth must come from sensitivity, outliers, activations or contexts. P6: error feedback or stochastic rounding for rollout weights, with a convergence theorem in place of \cref{cor:partial-sums}'s heuristic.
\end{enumerate}

% Background references listed in the bibliography but not cited in the text above.
\nocite{courbariaux2015binaryconnect,li2017training}

\bibliographystyle{plainnat}
\bibliography{refs}

\end{document}
